% Revised literature chain. Exactly three intervention and two speedup paragraphs.
% Source checks and interpretation limits: positioning.md and review reports.
\section{Related Work}
\label{sec:related-work}

\subsection{Sparsification interventions}
\label{sec:related-interventions}

Post-hoc thresholding reveals sparsity that an already learned
representation can tolerate. CATS thresholds small-magnitude activations
within gated FFNs, with both training-free and fine-tuning
settings~\citep{lee2024cats}. TEAL extends magnitude thresholding to attention
and FFN projections and uses block-wise greedy allocation to account for
their different sensitivities~\citep{liu2024teal}. Its placement and allocation
studies show why a uniform sparsity target need not be equally appropriate
throughout a model. These methods provide a reference for sparsity available
without retraining, while motivating the question of how training can change
that tolerance.

Training-time interventions allow the representation to adapt to sparse
computation. ReLU Strikes Back studies restoring sparse activations in
language models~\citep{mirzadeh2023relu}; ProSparse combines ReLU, progressive
$\ell_1$ pressure, and threshold shifting during continued
training~\citep{song2024prosparse}. ProSparse's schedule ablations show that
comparable FFN sparsity can carry different quality costs. Sparsing Law
examines how sparsity depends on regularization, data, and model
size~\citep{luo2025sparsing}, while \citet{cetin2026sparser} demonstrate highly
sparse FFNs induced by $\ell_1$ regularization during pretraining. Together,
these studies establish pressure as a way to reshape the quality--sparsity
trade-off. A local FFN zero fraction alone leaves a further accounting question:
how much of the complete model's computation do those zeros affect?

Extending sparsification beyond FFNs makes the choice of sites explicit.
Q-Sparse applies top-$k$ selection to attention and FFN projection inputs,
reports overall sparsity and activated-parameter budgets, and ablates both
the sparsifying nonlinearity and the straight-through
estimator~\citep{wang2024qsparse}. Spark Transformer combines sparse FFNs
with sparse attention through token-position selection and examines the role
of its low-cost predictor~\citep{you2025spark}.
Projection-input sparsity and attention-position sparsity affect different
work; our intervention ladder also considers feature coordinates in the
internal attention operands. Building on the existing
component studies, we examine how the additional effect of activation pressure
depends on thresholding and placement. Relating those effects to model-wide
product counts and topology-dependent reach connects the quality cost of an
intervention to the computation its zeros can affect.

\subsection{Inference speedup}
\label{sec:related-speedup}

Converting activation sparsity into speedup requires an execution strategy
suited to the workload. TEAL accelerates batch-one autoregressive decoding,
where skipping weight accesses is valuable, and reports up to $1.8\times$
speedup at 50\% sparsity under its projection-based
accounting~\citep{liu2024teal}.
\citet{cetin2026sparser} address batched matrix multiplication: their TwELL
format packs sparse activations as a projection produces them and supports
fused FFN computation, reducing the overhead of materializing and consuming
the sparse representation. They evaluate throughput for full-sequence forward
passes, a different workload from token-by-token decoding. Spark likewise couples its
sparsity mechanism with an implementation that reports decoding
gains~\citep{you2025spark}. These results show why the location and execution
of sparsity matter together. Skipped products and memory traffic must repay
indexing, packing, and launch overheads, while operations that remain dense
limit the gain of the complete model. A sparsity percentage alone does not
specify this balance.

This dependence makes kernel adaptation part of evaluating whether a
sparsification method transfers. With larger batches, different token masks
can reduce the weight-loading savings available from sparsity, as Spark's
batching analysis illustrates~\citep{you2025spark}. Block Q-Sparse introduces
structure into the mask to accommodate batched
execution~\citep{wang2024qsparse}; TwELL instead makes packing and fusion
compatible with tiled computation~\citep{cetin2026sparser}. Tensor shape,
sparsity structure, and data layout therefore matter alongside the zero
fraction. Iterative kernel generation with execution and profiling feedback
offers an emerging tool for exploring such
adaptations~\citep{ouyang2025kernelbench,luo2026agentic}.
Establishing a benefit still requires numerical correctness and complete-model
timings that include selection and conversion costs on the target workload.
These measurements test whether intervention-induced zero-product
opportunities translate into reduced execution time after sparse-execution
overheads are included.
